% !TeX root = ../../Book5.tex
% 纲要标题：### 27.6 陈特征、Todd 类与 Euler 示性数
% 纲要位置：几何学计划纲要.md，第 6172 行。
% 正文状态：未撰写；本文件目前是框架。

% 纲要内容开始
% 
% 以域上的纯维光滑射影概形及向量丛为计算主范围，使用有理 Chow 类；整数 Euler 数与含分母的特征类之间的关系需要核验。
% 
% * 向量丛的 Grothendieck 群与凝聚层的加性关系
% * 分裂原理下的陈特征及可加、可乘性
% * Todd 幂级数、低次数项及正规化
% * 光滑射影概形上的 Hirzebruch–Riemann–Roch 公式
% 
%   准确引用 $\chi(X,E)=\deg\bigl([\operatorname{ch}(E)\operatorname{td}(T_{X/k})]_{\dim X}\bigr)$；前面两种低维 Riemann–Roch 的证明不反向依赖本节。[R11]
% 
% * 曲线、曲面及射影空间的完整展开
% * 与代数曲线、曲面 Riemann–Roch 的独立核对
% 
% 纲要内容结束

\section{陈特征、Todd 类与 Euler 示性数}
\label{g5:sec:2706}

待撰写

\subsection{向量丛的 Grothendieck 群与凝聚层的加性关系}
\label{g5:subsec:2706:01}
待撰写

\subsection{分裂原理下的陈特征及可加、可乘性}
\label{g5:subsec:2706:02}
待撰写

\subsection{Todd 幂级数、低次数项及正规化}
\label{g5:subsec:2706:03}
待撰写

\subsection{光滑射影概形上的 Hirzebruch–Riemann–Roch 公式}
\label{g5:subsec:2706:04}
待撰写

\subsection{曲线、曲面及射影空间的完整展开}
\label{g5:subsec:2706:05}
待撰写

\subsection{与代数曲线、曲面 Riemann–Roch 的独立核对}
\label{g5:subsec:2706:06}
待撰写
% 纲要条目排为编号小节标题，并进入目录与书签；说明段仍保留于纲要注释中。
